\pdfinfoomitdate=1
\pdftrailerid{}
\pdfsuppressptexinfo=15
\documentclass[10pt,journal,compsoc]{IEEEtran}
\usepackage[T1]{fontenc}
\usepackage{amsmath,amssymb,booktabs,array,graphicx,cite,url}
\usepackage{algorithm,algorithmic}
\usepackage[nopatch=eqnum]{microtype}
\usepackage[hidelinks]{hyperref}
\hypersetup{pdftitle={FLOW-DC: A Large-Scale Distributed Downloader Utilizing Policy-Aware Adaptive Request Control},pdfauthor={Zi Deng, Nirav Merchant, Jeffrey J. Rodriguez}}
\newcommand{\grad}{PAARC-Grad}
\newcommand{\ratio}{PAARC-Ratio}
\makeatletter
\newcommand{\tableinput}[1]{\@@input #1 }
\makeatother
\newcommand{\clip}{\operatorname{clip}}
\newcommand{\floor}[1]{\left\lfloor #1\right\rfloor}
\newcommand{\ceil}[1]{\left\lceil #1\right\rceil}
\hyphenation{data-set data-sets Task-Vine}
\title{FLOW-DC: A Large-Scale Distributed Downloader Utilizing Policy-Aware Adaptive Request Control}
\author{Zi~Deng, Nirav~Merchant, and Jeffrey~J.~Rodriguez%
\IEEEcompsocitemizethanks{\IEEEcompsocthanksitem The authors are with the University of Arizona, Tucson, AZ, USA.}}
\markboth{FLOW-DC: Policy-Aware Adaptive Request Control}{Deng et al.: FLOW-DC}
\IEEEtitleabstractindextext{%
\begin{abstract}
Advances in deep learning depend on increasingly large datasets, yet assembling web-scale image collections remains a practical bottleneck for research groups with limited storage, bandwidth, and orchestration resources. URL-indexed collections shift acquisition to the user, while aggressive requests can trigger throttling and interrupt delivery. We present FLOW-DC, a Flexible Large-scale Orchestrated Workflow for Data Collection that combines asynchronous transfer, adaptive per-host request control, and distributed task execution. Its Policy-Aware Adaptive Request Control framework provides two complementary methods: PAARC-Ratio adjusts concurrency according to a baseline-relative delay budget, whereas PAARC-Grad also responds to persistent growth in excess application delay. Inspired by the probing and delay-reference principles of Google's Bottleneck Bandwidth and Round-trip propagation time control, both methods discover an operating region from observable HTTP behavior and integrate overload handling with accountable output generation. A TaskVine manager partitions manifests and coordinates workers without requiring a shared filesystem or synchronized arrival. Retained TaiBIF acquisitions deliver \TaiBIFRatioGoodput{} and \TaiBIFGradGoodput{} MB/s of verified end-to-end goodput for PAARC-Ratio and PAARC-Grad in one matched validation block. A separate four-block Open Images comparison achieves complete output coverage for both methods and img2dataset, while exposing the startup cost of adaptation on short jobs. Applications in biodiversity dataset construction, insect recognition, and weed identification demonstrate the practical relevance of the acquisition workflow. FLOW-DC provides a reusable path from distributed image URLs to verifiable, training-ready artifacts while making request control an explicit part of dataset construction.
\end{abstract}
\begin{IEEEkeywords}Dataset acquisition, adaptive concurrency, asynchronous HTTP, distributed workflows, biodiversity informatics.\end{IEEEkeywords}}
\input{sections/numbers}
\begin{document}
\maketitle
\IEEEdisplaynontitleabstractindextext
\IEEEraisesectionheading{\section{Introduction}\label{sec:intro}}
\IEEEPARstart{S}{tarting} from the breakthrough of AlexNet in large-scale image recognition~\cite{AlexNet}, deep learning has transformed classification, detection, and representation learning across scientific disciplines~\cite{DL-Survey}. Convolutional networks and vision transformers now provide complementary design choices to optimize performance for visual tasks~\cite{ViT,ConvNeXt}. To enhance effectiveness, developments ranging from BERT and generative pretraining in language~\cite{bert,gpt} to Swin transformers and YOLO in vision~\cite{swin,yolo} have utilized and depended on substantial training corpora. Datasets like ImageNet~\cite{imagenet}, Common Crawl~\cite{commoncrawl}, and OpenWebText~\cite{openwebtext} illustrate the breadth of this dependency. Acquiring a usable dataset is therefore a substantial part of the computational research process, rather than a negligible preliminary step.

The largest industrial laboratories can combine extensive storage, dedicated networking, and managed data-processing services. Academic groups more often assemble resources across institutional clusters, research clouds, and workstations. Even when training compute is available, data acquisition can consume a substantial portion of a project's preparation time. Provisioning more workers does not automatically solve this problem: data must be partitioned, transferred, checked, packaged, and made accessible to subsequent analysis. A practical system should support this workflow without assuming that every worker belongs to the same administrative environment.

Despite the breadth of accessible image repositories, converting URL-indexed collections into machine-learning-ready datasets remains demanding. iNaturalist, biodiversity portals, and curated web-image corpora distribute metadata that identifies externally hosted media assets. Users must handle various challenges such as changing availability, failed requests, duplicate rows, naming collisions, and storage layout. Metadata distribution is efficient for the publisher, but it transfers the work of acquisition and verification to each research group. The rate at which usable images become available directly affects the rate of research progress and model training.

Increasing request concurrency can improve utilization when individual transfers spend time waiting for remote responses. However, increasing it indiscriminately can also produce longer delays, HTTP refusals, and repeated work. Different public host services all differ in their capacity, policies, caching, and background demand. An HTTP error or delay increase is observable; the hidden mechanism that caused it usually is not. A useful acquisition tool should therefore treat these signals conservatively, respect explicit retry instructions, and maintain a record of what was actually delivered. It should not equate a high request count with a high rate of useful output.

Existing tools such as img2dataset provide effective URL-to-dataset pipelines, including parallel transfer, several output formats, configurable retries, and distributed execution~\cite{img2dataset}. However, these tools rely on configured process and thread counts rather than feedback-driven, per-host request admission, leaving concurrency tuning to the operator as source conditions change. FLOW-DC addresses these gaps by integrating per-host adaptive request admission, row-level outcome accounting, and artifact verification within an acquisition workflow designed for heterogeneous download resources.

\subsection{Automatic Concurrency Discovery under Unknown Server Conditions}
For a new source, an operator may know neither a useful concurrency setting nor how long that setting will remain appropriate. The best setting for one service or network path need not transfer to another. A conservative choice may leave capacity unused, whereas an aggressive choice can increase delay or provoke refusal. Repeating a manual search is both time-consuming and inconvenient, particularly when datasets span many hosts or when the available resources change between runs.

FLOW-DC makes this operating-point discovery part of acquisition itself. It begins with a small number of outstanding requests, observes application delay and overload feedback, and adjusts subsequent admission. This design can respond adaptively within a download job and can discover a new operating region when a previously useful baseline becomes stale. The practical value in this is that it reduces dependence on prior host knowledge and repeated operator intervention.

We take inspiration from Google's Bottleneck Bandwidth and Round-trip propagation time (BBR) congestion control~\cite{bbr}. BBR separates capacity probing from maintaining a low-delay reference. FLOW-DC carries these principles to HTTP acquisition, where the controllable quantity is the number of admitted requests rather than packet pacing. We introduce Policy-Aware Adaptive Request Control (PAARC), which combines a baseline-relative delay budget with the direction and persistence of excess-delay change to regulate request admission. PAARC requires neither access to server internals nor modifications to the transport stack.

FLOW-DC follows a manager--worker paradigm through TaskVine~\cite{taskvine}. A central manager partitions a dataset manifest and assigns acquisition tasks to workers. Each worker downloads asynchronously, applies the PAARC congestion rule, records row-level outcomes, and prepares verifiable output artifacts. Workers are easily able to join an ongoing workflow when their network access is compatible. This permits collaboration across resource pools without requiring a common filesystem or simultaneous startup.

The principal contributions are:
\begin{itemize}
\item An acquisition architecture that combines manifest partitioning, elastic task execution, asynchronous HTTP transfer, and verifiable output in one workflow.
\item An application-layer adaptive request controller, \paarc{}, that combines bounded exploration, a normalized excess-delay rate, and overload safeguards to discover useful concurrency settings.
\item A practical evaluation centered on verified goodput, completion time, and original-row coverage, together with applications in biodiversity and agricultural machine learning.
\end{itemize}
Section~\ref{sec:background} introduces the data and networking concepts. Section~\ref{sec:architecture} describes the architecture and controllers. Section~\ref{sec:results} presents the evaluation and discussion, and Section~\ref{sec:applications} describes research applications.

\section{Background}\label{sec:background}
\subsection{Image Dataset Primitives}
In modern machine-learning frameworks, data-loading utilities separate sample access from batching and scheduling. In PyTorch, the \texttt{Dataset} class defines how to obtain a sample, while a \texttt{DataLoader} organizes iteration, batching, and parallel loading~\cite{pytorch}. Downloading is an earlier concern: the referenced bytes must first be acquired, associated with their metadata, and stored in a form that these utilities can consume. These stages should remain distinct so that a training pipeline can choose its own decoding, resizing, and augmentation policy.

Notably, the scale of image datasets has grown substantially in the modern era. CIFAR-10 contains 60,000 small images~\cite{cifar10}; ImageNet (1.2 million images) and ImageNet-21K (11 million images) support much larger classification tasks~\cite{imagenet,imagenet21k}. Finally, web-derived collections such as LAION-5B~\cite{laion5b} distribute over 5 billion images for training in state of the art models. At this scale, distributing a manifest is often more practical than redistributing every image. However, a URL is a reference to a resource, not a guarantee that the resource remains available or unchanged. Dataset construction must preserve the distinction between the requested population and the subset successfully acquired.

Two storage patterns are especially useful for machine learning training. A directory of individual images is convenient for inspection and small datasets, but very large file counts increase filesystem metadata overhead. Sharded tar archives provide a sequential representation that packages images with associated metadata and can support streaming training~\cite{tar,webdataset}. FLOW-DC separates output layout from concurrency control so that the same acquisition mechanism can serve different downstream conventions. In the evaluation here, original hosted bytes are preserved; image validity and archive integrity are checked separately.

\subsection{Parallel HTTP Data Transfer From Manifest}
A dataset manifest is a tabular description of remote objects and their associated metadata. At minimum, it identifies the URL for each requested item. Additional columns may contain labels, taxonomic identifiers, or source attribution. Manifest preparation determines the scientific population; acquisition determines which members can be delivered. If invalid rows were silently discarded, success would be measured against a smaller set than was requested, overstating completeness. FLOW-DC therefore keeps every original row identifiable through to the output and records an explicit disposition for each.

HTTP defines request and response semantics~\cite{rfc9110}. Its HTTP/1.1 mapping operates over a reliable transport~\cite{http}, commonly TCP, whose original specification~\cite{tcp} has been consolidated in a current standard~\cite{tcp-current}. A client issues a GET request, receives headers and a status, and reads the response body. A successful status alone does not establish that the saved file is usable: the body may be incomplete, invalid, or incorrectly associated with a manifest row. Transfer, persistence, and verification therefore contribute separately to reliable acquisition. Concurrent downloads overlap waiting time across requests. FLOW-DC uses Python asynchronous I/O: coroutines yield while waiting, and an event loop schedules other ready operations. Coroutines are a general mechanism for cooperative execution~\cite{async}; the HTTP implementation uses \texttt{aiohttp}~\cite{aiohttp}. This avoids dedicating a Python thread to every pending network operation.

The receiver exposes several useful signals. HTTP 429 explicitly indicates too many requests and can carry a \texttt{Retry-After} instruction~\cite{http-codes}. HTTP 503 can also include retry guidance~\cite{rfc9110}. Other server errors, connection failures, or rising delay may justify reducing demand, but they do not uniquely identify a rate limiter. In particular, HTTP 403 is not, on its own, proof that a client has been blocked for excessive concurrency. FLOW-DC combines conservative admission rules with outcome records rather than attempting to infer an undocumented quota from a single response.

Request concurrency, request-start rate, and byte throughput are different quantities. Large objects may keep a request active for a long time without a high start rate; small objects may do the opposite. Likewise, a client can maintain many open requests while delivering little verified output. We use concurrency to control admission and verified bytes per end-to-end second to assess useful delivery.

\subsection{BBR Inspiration}
FLOW-DC utilizes a congestion control algorithm known as PAARC (Policy-Aware Adaptive Request Control). PAARC draws inspiration mainly on Google's BBR (Bottleneck Bandwidth and Round-trip propagation time)~\cite{bbr} algorithm, which tries to run at full bandwidth before queues fill rather than waiting for packet loss to signal congestion. BBR models a transport path using an estimated bottleneck delivery bandwidth, $B$, and a propagation-delay reference, $R$~\cite{bbr}. Their product expresses an inflight-data scale:
\begin{equation}
\mathrm{BDP}=BR,\qquad r_{\mathrm{pace}}=g_pB,\qquad W=g_wBR,
\label{eq:bbr}
\end{equation}
where $g_p$ and $g_w$ are pacing and window gains. With bandwidth in bytes per second and delay in seconds, BDP and $W$ are byte quantities. The original BBR design separates startup, draining, bandwidth probing, and propagation-delay refresh, as summarized in Fig.~\ref{fig:bbr}.

\begin{figure*}[t]\centering
\includegraphics[width=.96\textwidth]{figures/bbr-phases.pdf}
\caption{Conceptual phases of the original BBR design~\cite{bbr}.}
\label{fig:bbr}\end{figure*}

PAARC adopts the broader principles of controlled exploration, a low-delay reference, and feedback-driven regulation. Its application-layer observations and actions are different. The measured signal is TTFB (Time To First Byte), the interval from dispatching request headers to receiving the first body byte on the final HTTP hop. It includes effects of network transit, server processing, and client scheduling; it is neither packet RTT nor a direct measurement of the server queue. The control variable is an integer request limit. PAARC does not derive its limit from a packet-level BDP. Delay-change signals also have established precedent in transport research, including TIMELY~\cite{timely}. Our contribution is the construction and integration of application-delay feedback with request admission, public-host safeguards, distributed execution, and verifiable dataset output. Table~\ref{tab:bbr} makes the intended correspondence explicit.

\begin{table*}[t]\centering\small
\caption{BBR inspiration and application-layer realization in PAARC}
\label{tab:bbr}
\begin{tabular}{p{.18\textwidth}p{.32\textwidth}p{.42\textwidth}}\toprule
Design element & Original BBR & PAARC and PAARC-Ratio\\\midrule
Controlled quantity & Packet pacing and inflight bytes & Per-host integer request-admission limit\\
Delay reference & Minimum packet round-trip estimate & Low-percentile application-delay baseline\\
Exploration & Probe transport delivery capacity & Grow a utilized limit during startup\\
Feedback & Delivery-rate and transport-path observations & Successful-response delay and explicit overload feedback\\
Refresh & Temporarily reduce inflight data & Reduce admission, drain requests, collect fresh samples\\
Regulation & Transport-model gains & Delay-ratio or persistent excess-delay-rate rule\\
Output objective & Efficient transport delivery & Verified original-row bytes per end-to-end second\\\bottomrule
\end{tabular}\end{table*}

\section{Architecture Overview}\label{sec:architecture}
FLOW-DC organizes data acquisition into three stages: input preparation, coordination, and execution. A dataset manifest defines the requested objects; a manager converts partitions into tasks; workers acquire the objects and return their outputs and accounting records. Figure~\ref{fig:overview} presents this organization.

\begin{figure*}[t]\centering
\includegraphics[width=.97\textwidth]{figures/architecture.pdf}
\caption{Overview of FLOW-DC. Manifest partitions and configuration flow from the manager to workers. Workers apply PAARC to public-host requests and produce verifiable archives and row-level records.}
\label{fig:overview}\end{figure*}

\subsection{TaskVine}
TaskVine provides the manager--worker execution layer~\cite{taskvine}. The manager declares a task's command, resource requirements, input files, and output files. Workers connect to the manager, receive runnable tasks and their declared inputs, and return declared outputs. This design avoids requiring a shared filesystem between all participants. It also permits workers to join after task submission, provided that the credentials, and network connectivity are compatible.

\begin{figure}[t]\centering
\includegraphics[width=.97\columnwidth]{figures/taskvine.pdf}
\caption{TaskVine manager--worker organization. The scheduler transfers declared inputs and outputs and assigns available tasks to connected workers.}
\label{fig:taskvine}\end{figure}

FLOW-DC concentrates on the acquisition stage: partitions have similar execution structure, but their duration depends on remote hosts and their objects. TaskVine supplies scheduling and file movement; FLOW-DC supplies manifest semantics, the HTTP client, request control, and output verification. A scheduler's successful task return is not treated as sufficient evidence of successful acquisition. Additionally, ther workflow engines such as Snakemake~\cite{snakemake}, Nextflow~\cite{nextflow}, and Pegasus~\cite{pegasus} were considered and tested but determined to be less suitable.

\subsection{Policy-Aware Adaptive Request Control}
PAARC controls a per-host concurrency limit $C_t$. Outstanding requests can temporarily exceed a newly reduced limit because the controller does not cancel otherwise healthy transfers merely to enforce an instantaneous reduction. Further admissions wait until capacity is available. A configured upper bound remains necessary to protect finite client resources; automatic discovery does not imply unlimited concurrency.

\subsubsection{Measurement and baseline}
Let $x_i$ denote the application delay of an eligible successful HTTP response, measured in seconds from final-hop request-header dispatch to the first response-body byte. The response becomes eligible after body completion. Admission waiting, connection-pool waiting, and preceding redirect hops are excluded from this signal. Consequently, $x_i$ is useful feedback about delivered requests, but does not describe all waiting experienced by a manifest row.

At a control opportunity, at least five eligible observations are required. Both dispatch and observation must fall within a two-second freshness window, and consumed observations are not reused. Let $p_{10,t}$ and $p_{50,t}$ be their interpolated tenth percentile and median. The baseline $b_t$ is initialized from $\max(p_{10,t},\epsilon)$ and updated when a lower eligible percentile appears. The numerical floor is $\epsilon=10^{-6}$ s. A 120-s baseline lifetime triggers a refresh: admission falls to $C_{\min}$, outstanding work drains, and fresh observations dispatched after the probe begins establish the new reference. Refresh can therefore move the baseline upward as well as downward.

For an eligible update separated by $\Delta t$ seconds from the preceding eligible update, define excess delay and elapsed-time smoothing by
\begin{align}
q_t^{\rm raw}&=\max(0,p_{50,t}-b_t),\\
\alpha_q&=1-\exp(-\Delta t/\tau_q),\\
\bar q_t&=(1-\alpha_q)\bar q_{t-1}+\alpha_q q_t^{\rm raw},\\
\bar d_t&=(1-\alpha_q)\bar d_{t-1}+\alpha_q p_{50,t}.
\label{eq:delay}
\end{align}
Here $\bar q_t$ and $\bar d_t$ have units of seconds. The baseline-normalized excess-delay rate is estimated by the finite difference
\begin{align}
s_t&=\frac{\bar q_t-\bar q_{t-1}}{\Delta t\max(b_t,\epsilon)},\\
\alpha_s&=1-\exp(-\Delta t/\tau_s),\\
\bar s_t&=(1-\alpha_s)\bar s_{t-1}+\alpha_s s_t,
\label{eq:delay-rate}
\end{align}
with $s_t$ and $\bar s_t$ in s$^{-1}$. This is a sampled rate of change in excess delay relative to the baseline, rather than a partial derivative of an optimization objective. Both time constants are 1 s. The first eligible sample set initializes the delay state and holds the limit; the first subsequent finite difference initializes the smoothed rate directly. A baseline change or a five-second stale gap clears delay-rate history. These resets prevent a reference change from being mistaken for a queue trend.

\subsubsection{Exploration and common safeguards}
PAARC and its ratio alternative begin at $C_0=4$, with $C_{\min}=2$. Let $I_t$ be the number of outstanding requests. Growth requires utilization $I_t\geq0.8C_t$. During exploration, the update is
\begin{equation}
C_{t+1}=\clip\!\left(\min(C_t+64,\ceil{\gamma C_t}),C_{\min},C_{\max}\right).
\label{eq:startup}
\end{equation}
The growth factor $\gamma$ is reported with each evaluation configuration. Exploration ends when smoothed delay reaches the budget $D_t=\rho b_t$, with $\rho=2$ in the reported acquisitions. Three eligible updates with $\bar s_t\leq-0.1$ s$^{-1}$ reenable exploration, subject to the same delay-budget and utilization checks.

Explicit overload feedback takes precedence over growth or ordinary regulation. Connection failures, HTTP 408 or 429, and server errors contribute to this conservative signal; locally identified persistence errors are not interpreted as remote overload. The limit is halved, delay-rate state is cleared, and a one-second recovery interval is entered. The same hard decrease applies when $\bar q_t\geq\max(0.001\,\mathrm{s},4b_t)$. Positive growth is held during recovery. Sparse observations also cause a hold; lack of evidence is not interpreted as unused server capacity. Valid \texttt{Retry-After} instructions impose admission embargoes independently of the controller's desired limit.

\subsubsection{PAARC regulation}
Outside exploration, \paarc{} first reduces $C_t$ by a factor of 0.85 when $\bar d_t\geq D_t$. Otherwise, an excess delay of at least 1 ms activates feedback from the smoothed normalized excess-delay rate. A delay rate at or above 0.25 s$^{-1}$ for two consecutive eligible updates causes the same decrease; the first such update holds the limit. A delay rate at or above 0.1 s$^{-1}$ also holds growth. When these conditions are absent and utilization is sufficient, the limit increases by
\begin{equation}
\Delta C_t=\min\left(8,\max\left(1,\ceil{0.02C_t}\right)\right).
\label{eq:paarc-growth}
\end{equation}
Soft decreases clear delay-rate state and enter recovery. All decreases use floor rounding and all resulting limits are clipped. This rate supplies an early indication of increasing excess delay, while the absolute budget prevents a stable but elevated delay from being accepted solely because its rate of change is near zero. Neither signal identifies the remote cause uniquely.

\subsubsection{Alternative controller: PAARC-Ratio}
We also evaluate \ratio{}, an alternative controller that shares the measurement, exploration, and priority safeguards above. During regulation, \ratio{} maps the delay budget to a bounded multiplicative factor:
\begin{align}
r_t&=\clip\!\left(\sqrt{D_t/\bar d_t},0.85,1.10\right),\\
\widehat C_{t+1}&=\floor{C_t r_t+1}.
\label{eq:ratio}
\end{align}
The proposal is accepted when it decreases the limit or the current limit is utilized; otherwise admission remains unchanged. Final clipping enforces the configured bounds. The additive unit and integer rounding are part of the rule, so the equilibrium is not obtained by simply setting $\bar d_t=D_t$. Ordinary ratio decreases preserve the filtered state; hard safeguards still reset it. This provides continuous regulation around the selected delay allowance without requiring a transport-bandwidth estimate.

These rules form a priority-ordered transition system with exploration, probe, recovery, and observation state variables. The logical phases are illustrated in the supplement. They are not a claim that PAARC implements BBR's transport state machine. The supplement specifies branch order, parameter values, and source mappings needed to reproduce the decisions.

\subsection{Input}
FLOW-DC accepts a dataset manifest specifying items to download and their associated metadata. Rows correspond to original requested items, and columns describe URLs, labels, identifiers, or attribution. The maintained acquisition path consumes prepared Parquet partitions containing HTTP or HTTPS URLs. Dataset-specific preparation can convert other source metadata into this representation; support for an upstream metadata format does not imply support for another transfer protocol.

For example, an iNaturalist acquisition can derive URLs and taxonomic labels from its open-data metadata~\cite{inat,inat-open}. Preparation must retain the requested population, licenses, source identifiers, and sampling rule. Invalid URLs and duplicate rows retain explicit identities and dispositions. The research path accepts a restricted set of scalar metadata types and rejects unsupported types before acquisition. This contract makes metadata preservation testable without claiming unrestricted handling of every possible dataframe schema.

\subsection{Coordination}
The workflow prepares manifest partitions, stores them as Parquet files, and submits a download task for each partition. Partition size controls scheduling granularity: smaller partitions expose more parallel tasks, whereas larger partitions amortize startup and output-finalization overhead. It must be selected with object sizes, runtime, memory, and archive cost in mind rather than interpreted as a concurrency setting.

\textbf{Grouping by host.} The preparation utilities support sequential partitions and host-oriented grouping. In the latter, hosts are ordered by their URL counts and assigned to the currently least-loaded partition. If $c_h$ is a host's row count and $S_j$ is the load of partition $j$, the assignment is
\begin{equation}
j^*=\arg\min_j S_j,\quad A(h)=j^*,\quad S_{j^*}\leftarrow S_{j^*}+c_h.
\label{eq:partition}
\end{equation}
This retains host locality, but does not guarantee equal bytes or equal execution time: row counts are only a proxy for work. A very large host can also dominate a partition. Where work for one host spans concurrent tasks, independent local controllers do not constitute a global host limit. FLOW-DC's shared-admission path can coordinate a common authority when explicitly configured. Its leases and request accounting must cover the selected workers; merely adding workers does not establish that property.

\textbf{Task declaration.} Each task declares its partition, configuration, downloader, required Python modules, and outputs. The staging closure includes the HTTP, method-selection, and integrity modules used by the selected path. Workers return summaries and declared data artifacts according to the configured workflow. Configuration and source identities accompany experiment records, so a result can be associated with the code and settings that produced it. The scheduler and the scientific accounting layer retain separate responsibilities.

\subsection{Execution}
A worker receives declared inputs and executes the batch downloader in its task environment. It validates the partition, initializes the HTTP session and observation buffers, and runs asynchronous acquisition under per-host admission. DNS caching and connection reuse reduce repeated setup where supported by the session. Request outcomes feed the controller, while row dispositions and saved artifacts feed the integrity ledger.

\begin{figure}[t]\centering
\includegraphics[width=.96\columnwidth]{figures/worker.pdf}
\caption{Worker acquisition and output pipeline. Admission decisions use HTTP observations; credited output additionally requires persistence and independent verification.}
\label{fig:worker}\end{figure}

A successful transfer is persisted through the integrity path before it is credited as delivered output. The journal supports reconciliation after interrupted writes and preserves the distinction between original rows, acquisition attempts, and eligible artifacts. Archive verification checks membership and content against retained records. A task that stops during acquisition or finalization remains incomplete until its retained artifacts establish otherwise. This is particularly important when substantial time is spent packaging or checking data after HTTP activity has ended.

\subsection{Features}
\textbf{Elastic execution.} Task-based scheduling allows resources to participate without synchronized startup. Practical deployment still depends on runtime compatibility, reachable endpoints, and available storage. Distributed execution is a capability of the architecture; its benefit depends on whether the source, manager, and storage path have sufficient capacity.

\textbf{Output flexibility.} Acquisition and training representation remain separate. Image directories facilitate inspection, while sharded archives reduce the burden of managing large numbers of files. Retaining original hosted bytes allows downstream researchers to select their own decoding and transformation policies.

\textbf{Accountable recovery.} Resume and reconciliation use the existing journal and artifacts to distinguish completed work from work that must be attempted again. Repetition is recorded, rather than silently replacing an unfavorable first attempt. These mechanisms improve operational reliability without making incomplete output appear complete.

\textbf{Extensibility.} Dataset preparation, scheduling, HTTP acquisition, control decisions, and verification have distinct interfaces. A new dataset usually changes manifest construction rather than controller logic. A new transfer mechanism would require its own measurement and integrity qualification; it cannot inherit HTTP-specific claims solely by sharing the scheduler.

\section{Experimental Results and Discussion}\label{sec:results}
The evaluation emphasizes practical acquisition: how quickly requested, verifiable bytes become available to downstream processing. We compare \ratio{} and \grad{} with img2dataset, and examine the relationship between transfer activity and end-to-end completion. The measured acquisitions use public image services. No generated origin or artificial image payload supplies the measured results.

\subsection{Evaluation Design and Analysis}
\textbf{Datasets and execution.} The Open Images acquisition uses public images hosted by the Common Visual Data Foundation (CVDF)~\cite{cvdf}. Four independently initialized evaluation blocks each contain 512 frozen URLs. Methods receive the same manifest within a block, with approximately 159--165 MB of hosted image bytes. The longer TaiBIF acquisition uses a frozen 32,768-row moth-image manifest from Taiwan's biodiversity infrastructure~\cite{taibif}. Both studies execute on Jetstream2 guests. Each acquisition uses one downloader guest; manager and origin-role guests support operations and do not serve experimental images. Source snapshots, settings, input hashes, and available environment records are identified in the supplement.

Public hosts are observational environments. Their background demand and caches are outside the client's control, and a given host need not expose a stable concurrency threshold. Requests preserve the selected URL population and record failures. Within the Open Images evaluation, method order is randomized inside blocks and settings are selected on disjoint inputs before evaluation. These small samples characterize the retained jobs; they are not a population-wide performance guarantee.

\textbf{Matched tool configuration.} img2dataset 1.47.0~\cite{img2dataset} uses one process and 192 threads for the Open Images comparison. Resizing and all reencoding are disabled, and retries are disabled for this bounded comparison. Both PAARC methods begin at four requests with a limit of 192, growth factor 1.5, delay ratio 2, and baseline lifetime 120 s. The output verifier applies a common original-row and byte-crediting rule. The clients retain their native connection and timeout implementations, which are documented rather than assumed equivalent. The comparison concerns these complete acquisition configurations, not a controlled substitution of one controller inside an otherwise identical HTTP client.

\textbf{Primary measures.} Let $\mathcal R$ be the original set of requested rows, $z_i\in\{0,1\}$ indicate independently verified eligible output for row $i$, and $B_i$ be that output's byte count. We report
\begin{equation}
G=\frac{\sum_{i\in\mathcal R}z_iB_i}{t_{\rm verified}-t_{\rm start}},
\qquad
Q=\frac{\sum_{i\in\mathcal R}z_i}{|\mathcal R|}.
\label{eq:goodput}
\end{equation}
$G$ is end-to-end goodput and $Q$ is original-row coverage. A duplicate manifest row retains its own disposition; byte credit follows the verifier's row contract, not the number of filesystem objects. The job clock includes the per-job execution and output verification boundary. Environment installation, batch orchestration, final collection, and VM cleanup are separate costs. All rates in this paper use decimal MB/s, with $1\,\mathrm{MB}=10^6$ bytes.

We also report the HTTP-active span, native-process time, output-processing time, and observed concurrency. Time-average concurrency integrates activity over the stated interval; a peak is only the maximum observed overlap. Neither is interchangeable with the admission limit. An active interval begins at different native instrumentation points in FLOW-DC and img2dataset, so cross-tool concurrency counts require caution even when their end-to-end byte accounting agrees.

\textbf{Analysis.} Runs, rather than requests, are the experimental units. For paired Open Images blocks, log goodput ratios compare each PAARC method with img2dataset. Their exponentiated mean and ordinary Student-$t$ interval summarize the retained blocks. With only four blocks, these are exploratory, unadjusted intervals. Interrupted and unstarted jobs remain in the inventory and are not imputed as complete observations.

\subsection{Single-Machine Acquisition and Tool Comparison}
Fig.~\ref{fig:cvdf} summarizes Open Images acquisition across the evaluated configurations. In the four matched blocks, img2dataset at 192 threads and both PAARC methods deliver complete verified coverage. img2dataset averages \CVDFToolGoodput{} MB/s and \CVDFToolTime{} s, while \ratio{} and \grad{} average \CVDFRatioGoodput{} and \CVDFGradGoodput{} MB/s, with completion times of \CVDFRatioTime{} and \CVDFGradTime{} s. Thus, the 192-thread img2dataset configuration is faster for this small, readily served workload. The geometric goodput ratios relative to this matched reference are 0.875 (95\% interval 0.771--0.993) and 0.858 (0.778--0.946), respectively.

% Temporarily hidden until the additional Q, Active and Time values are supplied.
% To restore, uncomment this table block and remove the reserved counter below.
% \begin{table}[t]\centering\footnotesize
% \setlength{\tabcolsep}{3pt}
% \caption{Open Images acquisition by configuration}
% \label{tab:cvdf}
% \begin{tabular}{lrrrrr}\toprule
% Method & Limit & MB/s & Time (s) & $Q$ (\%) & Active$^\dagger$\\\midrule
% \tableinput{sections/cvdf-table.tex}
% \bottomrule\end{tabular}
% \par\vspace{3pt}\raggedright\footnotesize Limit is configured concurrency; available means summarize four runs. N/A denotes unavailable measurements, including completion metrics for all five stopped 384-limit runs. $^\dagger$Mean of native time-average activity. $^\ddagger$Approximately two simultaneous requests observed, not a measured mean or requests/s.
% \end{table}
\stepcounter{table} % Reserve Table 2 so subsequent tables keep their numbering.
\begin{figure*}[t]\centering
\includegraphics[width=.96\textwidth]{figures/cvdf-comparison.pdf}
\caption{Open Images goodput by configuration: bars show means and points show individual runs; N/A marks five stopped runs, not zero goodput. The 96- and 384-thread tests used a different machine and are excluded from the matched statistical analysis.}
\label{fig:cvdf}\end{figure*}

The result is informative for automatic discovery. These jobs last only about 13--15 s end to end, and neither adaptive method reaches a long-running regime. Mean native activity is near seven requests for PAARC, while the tool starts with a much larger thread pool. Independent verification takes approximately 3.4--3.6 s per job. Startup and fixed per-job work therefore consume a substantial fraction of the useful acquisition interval. The 120-s baseline-refresh mechanism is not exercised by these short jobs. Complete output coverage demonstrates accountable delivery; it does not demonstrate the benefit of every control mechanism.

\textbf{Configuration sensitivity.} At concurrency 96, the four verified end-to-end goodputs are 6.12, 6.33, 6.74, and 6.26 MB/s (mean 6.36 MB/s). At concurrency 384, all five runs were stopped for time reasons after severe slowdown, with approximately two simultaneous requests observed. These runs did not complete the acquisition objective; their completed-run goodput and completion time are N/A.

\subsection{Response Analysis and End-to-End Cost}
The longer TaiBIF validation separates network activity from the full acquisition pipeline. One matched block has complete retained processing records for both methods. Here, complete means that the bounded job and its output accounting finished; it does not mean that every row in the manifest was downloaded. \ratio{} verifies 26,969 of 32,768 rows (82.30\%), and \grad{} verifies 28,204 (86.07\%). Their verified output is 4.143 and 4.341 GB, respectively.

Figure~\ref{fig:timing} shows why HTTP-period throughput alone can overstate the rate at which a usable dataset becomes available. For \ratio{}, 181.05 s of HTTP activity sits within 284.11 s of native execution and 444.60 s of end-to-end processing. The corresponding \grad{} intervals are 181.19, 287.17, and 450.43 s. Dividing verified bytes by the HTTP span gives 22.88 and 23.96 MB/s; the end-to-end rates are \TaiBIFRatioGoodput{} and \TaiBIFGradGoodput{} MB/s. The additional time includes native work outside the HTTP-active span and postprocessing or verification. The retained timing does not attribute every residual second to a single subsystem.

\begin{figure*}[t]\centering
\includegraphics[width=.97\textwidth]{figures/taibif-timing.pdf}
\caption{Measured TaiBIF validation block 2. Left: HTTP span, other native time, and postprocessing. Right: the same verified-byte numerator divided by the HTTP span or the complete job interval. The two arms acquire different amounts within a bounded window.}
\label{fig:timing}\end{figure*}

The full-HTTP-span mean activity is 53.06 requests for \ratio{} and 54.55 for \grad{}, with peaks of 131 and 145. In the designated 30--90 s window, both have median activity 60. These measurements show useful adaptive acquisition at substantially lower activity than their configured ceilings. They do not establish sustained activity above 150, an optimal operating point, or a server-side throttling threshold. A larger limit is permission to admit work, not evidence that such work is present or beneficial.

Across the twelve retained validation dispositions, three jobs complete, two stop on processing deadlines, and seven have no completed validation execution. Only the block above provides both complete arms. The additional complete \ratio{} job achieves 8.65 MB/s, while the stopped and unstarted records remain visible in the supplement. There is no matched img2dataset arm in this TaiBIF validation, so the Open Images tool comparison cannot be transferred to this host. The central positive result is practical verified acquisition and interpretable accounting, rather than a replicated superiority claim.

\subsection{Distributed Execution Qualification and Workload Sizing}
We additionally qualify the matched TaskVine execution path on iNaturalist Open Data and CVDF Open Images. These short jobs serve a different purpose from the preceding acquisition comparisons: they check partition placement, returned artifacts, and the resource requirements for a later long-run study. They are not independent efficacy replications, and their measured durations are not substituted for a ten-minute evaluation.

The manager has 16 virtual CPUs and 60 GiB of memory, while each downloader has eight virtual CPUs and 30 GiB. Each task receives eight cores and a 24-GiB memory budget. One- and two-worker iNaturalist jobs and one-worker Open Images jobs use 512 and 2,048 original rows. The two PAARC settings are unchanged. The tool additionally exercises the predefined process/thread configurations $(1,32)$, $(1,192)$, and $(8,24)$. Two-worker jobs halve the per-worker request ceiling or thread count, maintaining an aggregate ceiling of 192. Each worker runs its own local controller; a shared global adaptive authority is not evaluated. Every credited partition is independently verified after return to the manager.

\input{sections/engineering-status}

\begin{figure*}[t]\centering
\includegraphics[width=.97\textwidth]{figures/engineering-duration.pdf}
\caption{Measured 2,048-row execution qualification. Each bar is one completed engineering configuration, not a mean of independent efficacy replications. Color shows the HTTP span; gray shows the remainder of the manager interval. The two-worker span is the maximum individual worker span.}
\label{fig:engineering}\end{figure*}

\begin{table*}[t]\centering\small
\caption{Engineering measurements used for workload sizing}
\label{tab:engineering}
\begin{tabular}{lr lrr@{\hspace{2em}}rr}\toprule
Source & Workers & Method (processes, threads) & HTTP 512 (s) & HTTP 2,048 (s) & Total 2,048 (s) & Coverage (\%)\\\midrule
\tableinput{sections/engineering-table.tex}
\bottomrule\end{tabular}
\par\vspace{3pt}\footnotesize One completed engineering attempt per configuration and size; missing measurements are shown as dashes. Coverage refers to the 2,048-row job. Original stopped attempts remain in the supplement and artifact inventory.
\end{table*}

These measurements also expose a workload-sizing constraint. A common manifest must contain enough work for the fastest configuration to acquire images for at least 600 s, while allowing every candidate to finish native processing and manager verification within finite deadlines. Let $\widehat H_m(n)$, $\widehat N_m(n)$, and $\widehat T_m(n)$ denote engineering estimates of HTTP, native-process, and manager elapsed time for method $m$ and $n$ original rows. The admission rule requires
\begin{align}
\min_m\widehat H_m(n)&\geq600\ \mathrm{s},\nonumber\\
1.25\widehat N_m(n)&<1800\ \mathrm{s},\nonumber\\
1.25\widehat T_m(n)+60\ \mathrm{s}&<2400\ \mathrm{s}
\quad\text{for every }m.
\label{eq:workload-admission}
\end{align}
Response-byte and storage limits apply in addition. The engineering model separates an estimated fixed cost from per-row work using the two sizes. This model is an admission heuristic, not a measured long-run performance law: adaptation, public-service variability, and storage behavior can change its slope. A refused workload therefore identifies an unresolved experimental requirement, not proof that longer acquisition is intrinsically impossible. Longer-run speed, causal mechanism effects, and general scalability require their own completed measurements.


\subsection{Discussion, Limitations, and Reproducibility}
Automatic concurrency discovery is most valuable when prior host knowledge is limited and acquisitions last long enough to respond to observations. It reduces an operator's dependence on a manually chosen setting and can respond when the environment changes. Its benefit is conditional, however. Short jobs can favor a well-provisioned tool immediately; noisy or sparse observations can delay growth; and a low-delay response does not guarantee that the client is below every provider policy threshold. PAARC responds to observable signals and explicit retry instructions, not to a privileged view of server policy.

The historical comparative evidence spans two public services and a limited number of independent jobs. Remote changes, order effects, cache state, and differing native-client behavior constrain generalization. Incomplete TaiBIF validation limits paired inference. The additional engineering checks qualify distribution and output accounting; they do not supply missing long-run experimental replication. Likewise, application adoption establishes utility of the acquisition workflow, not a retrospective evaluation of the current control rules.

Reproducibility has two levels. The accompanying analysis package rebuilds every figure and table from included measured summaries and conceptual diagram code, with a claim-to-evidence mapping and source hashes. Repeating acquisition against the same public URLs can encounter changed bytes or availability and therefore cannot promise identical results. Retaining manifests, first-attempt outcomes, verification rules, timing boundaries, and software identities allows those differences to be examined rather than hidden. This separation is essential for a downloader whose scientific output depends partly on infrastructure outside its authors' control.

\section{Applications}\label{sec:applications}
FLOW-DC grew from the practical need to acquire large image collections for biodiversity and agricultural machine learning. These projects connect public metadata, distributed image hosts, dataset curation, and model development. Earlier versions of the acquisition workflow supported the projects below, as reported by the contributing authors. The publications document the resulting datasets and learning tasks. The current \ratio{} and \grad{} rules are evaluated separately in Section~\ref{sec:results}.

\subsection{Insect Identification from Citizen Science}
InsectNet combines image acquisition, representation learning, classification, and deployment for insect identification~\cite{insectnet}. Its data pipeline draws on iNaturalist imagery, with approximately six million labeled images spanning 2,526 insect species. The earlier study, \emph{Deep learning powered real-time identification of insects using citizen science data}~\cite{insect-preprint}, belongs to the same research lineage. These are related stages of one application family, rather than two independent performance evaluations.

The published InsectNet pipeline calls its downloader iNaturalist Scalable Download (iNatSD). The authors identify this acquisition work as part of the lineage leading to FLOW-DC. Its practical requirements---large URL lists, taxonomic metadata, distributed execution, and reusable image output---motivate the interfaces described here.

\subsection{BioTrove and Arboretum}
BioTrove provides a large curated biodiversity image collection and supports models that connect visual and textual information~\cite{biotrove,clip}. The published dataset contains approximately 161.9 million images across 366.6 thousand species. The earlier \emph{Arboretum: A Large Multimodal Dataset Enabling AI for Biodiversity} manuscript~\cite{arboretum} is the first version of the same arXiv record and reports an earlier dataset scale. It is therefore treated as an earlier dataset iteration rather than a separate deployment.

At this scale, acquisition is a workflow problem as much as a transfer problem. Manifest partitions must retain taxonomic and attribution metadata, workers must produce manageable artifacts, and failures must remain visible for later curation. An earlier FLOW-DC workflow supported this acquisition effort. That experience supports the system's emphasis on partitioned scheduling and archive-oriented output, while the current row-level verification strengthens the accountability of subsequent acquisitions.

\subsection{Weakly Supervised Insect Detection}
\emph{Zero-shot insect detection via weak language supervision} studies how image and language information can support insect detection without conventional task-specific supervision~\cite{weak-insects}. The acquisition workflow contributes reusable image collections and metadata that can serve tasks beyond the original classification objective. This use case emphasizes why FLOW-DC preserves the relationship between source rows and delivered images: a dataset assembled once may later support a different sampling rule, annotation strategy, or model family. Downloading is one enabling component; the downstream detection method remains a separate scientific contribution.

\subsection{Weed Identification and Reusable Research Infrastructure}
WeedNet develops a global-to-local approach to weed species identification, with a global model spanning 1,593 species and regional specialization~\cite{weednet}. The work appeared first as a 2025 preprint and was published online in \emph{Nature Communications} in September 2026. The acquisition workflow supported its dataset construction. This application illustrates the need to reuse the same acquisition machinery across biological groups while changing manifest selection and downstream learning procedures.

Across these projects, the recurring requirement is a dependable transition from public image references to organized research data. Research cyberinfrastructure such as CyVerse provides complementary environments for managing and sharing scientific analyses~\cite{cyverse}. FLOW-DC addresses the acquisition portion of that ecosystem. Its value lies in making large transfers easier to organize, adapt, inspect, and reuse, while preserving the source metadata and operational evidence needed by subsequent curation and analysis.

\section{Conclusion}
FLOW-DC connects large-scale dataset construction to the practical constraints of academic computing and public image services. Its manager--worker architecture separates manifest preparation, scheduling, asynchronous acquisition, and verified output, while \ratio{} and \grad{} make per-host concurrency a continuously informed decision. The two methods share startup and acquisition safeguards but use different delay-feedback rules during regulation. Measured public-host results demonstrate useful delivery and complete short-job output accounting; they also show that adaptation has a startup cost. Distributed engineering checks extend output-accounting validation to two downloader guests. Longer matched acquisitions remain necessary to assess sustained goodput and scaling benefits. Together with its use in biodiversity and agricultural research, the system offers a practical foundation for acquiring reusable datasets without requiring users to build a downloader and orchestration stack for each project.

\begingroup
\emergencystretch=1em
\IEEEtriggeratref{37}
\bibliographystyle{IEEEtran}
\bibliography{references}
\endgroup
\end{document}
